\documentclass[tikz,border=12pt]{standalone}
\usepackage{fontspec}
\usepackage[mathrm=sym]{unicode-math}
\setmainfont{texgyrepagella-regular.otf}[
    BoldFont = texgyrepagella-bold.otf,
    ItalicFont = texgyrepagella-italic.otf,
    BoldItalicFont = texgyrepagella-bolditalic.otf
]
\setmathfont{texgyrepagella-math.otf}[AutoFakeBold = 2.5]
\usepackage{tikz}
\usetikzlibrary{arrows.meta, calc}

\begin{document}
\begin{tikzpicture}[>=Stealth, line cap=round, line join=round]

  % =============================================================
  % LEFT PANEL: E vs J (Two-band crossing)
  % =============================================================
  \begin{scope}[shift={(0,0)}]
    % Axes
    \draw[line width=0.9pt, ->] (0, 0) -- (4.5, 0) node[anchor=north west, font=\large] {$J$};
    \draw[line width=0.9pt, ->] (0, 0) -- (0, 4.2) node[anchor=south east, font=\large] {$E$};

    % Crossing point J_c
    \def\xc{2.35}
    \def\yc{1.25}
    \draw[line width=0.8pt, dotted] (\xc, 0) -- (\xc, \yc);
    \node[anchor=north, font=\normalsize] at (\xc, -0.05) {$J_c$};

    % Unperturbed diabatic dashed crossing segments near J_c
    \draw[line width=0.8pt, dashed] (\xc - 0.45, \yc - 0.3) -- (\xc + 0.45, \yc + 0.35);
    \draw[line width=0.8pt, dashed] (\xc - 0.45, \yc + 0.1) -- (\xc + 0.45, \yc - 0.1);

    % Lower branch (perturbed Yrast band): starts from (0,0), follows band 1 then switches smoothly to band 2
    \draw[line width=1.2pt]
      (0, 0) to[out=10, in=205] (1.5, 0.5)
      to[out=25, in=210] (\xc, \yc - 0.12)
      to[out=30, in=220] (3.1, 1.7)
      to[out=40, in=235] (3.8, 2.5);

    % Upper branch (perturbed Yrare band): starts from (0, 1.1), follows band 2 then switches to band 1
    \draw[line width=1.2pt]
      (0, 1.1) to[out=0, in=195] (1.5, 1.15)
      to[out=15, in=210] (\xc, \yc + 0.12)
      to[out=45, in=245] (3.1, 2.7)
      to[out=65, in=255] (3.5, 3.5);

    % Labels of the two bands (positioned alongside curves without clipping)
    \node[anchor=east, font=\large] at (3.0, 3.1) {$\mathcal{J}_1$};
    \node[anchor=west, font=\large] at (3.9, 2.4) {$\mathcal{J}_2$};
  \end{scope}

  % =============================================================
  % RIGHT PANEL: \mathcal{J} vs \omega^2 (Backbending S-curve)
  % =============================================================
  \begin{scope}[shift={(5.6,0)}]
    % Axes
    \draw[line width=0.9pt, ->] (0, 0) -- (4.5, 0) node[anchor=north west, font=\large] {$\omega^2$};
    \draw[line width=0.9pt, ->] (0, 0) -- (0, 4.2) node[anchor=south east, font=\large] {$\mathcal{J}$};

    % Levels J1 and J2
    \def\yJone{1.2}
    \def\yJtwo{3.0}

    \draw[line width=0.8pt] (-0.08, \yJone) -- (0.08, \yJone);
    \node[anchor=east, font=\large] at (-0.12, \yJone) {$\mathcal{J}_1$};

    \draw[line width=0.8pt] (-0.08, \yJtwo) -- (0.08, \yJtwo);
    \node[anchor=east, font=\large] at (-0.12, \yJtwo) {$\mathcal{J}_2$};

    % Dashed asymptotic lines
    \draw[line width=0.7pt, dashed] (0, \yJone) -- (3.6, \yJone);
    \draw[line width=0.7pt, dashed] (0, \yJtwo) -- (3.6, \yJtwo);

    % Dash-dotted tangent line through the backbending region
    \draw[line width=0.8pt, dash pattern=on 5pt off 2pt on 1.5pt off 2pt]
      (1.3, \yJtwo) -- (2.9, \yJone);

    % S-shaped backbending curve
    \draw[line width=1.3pt]
      (0, \yJone) -- (1.8, \yJone)
      to[out=0, in=-90] (2.3, 1.6)
      to[out=90, in=-35] (1.85, 2.5)
      to[out=145, in=180] (2.2, \yJtwo)
      -- (3.6, \yJtwo);
  \end{scope}

\end{tikzpicture}
\end{document}
